\section{Structural limits}\label{sec:limits}

This section records what goes wrong when one of the ingredients is
removed. Each item is either a proved statement or a cited result; none is
a lower bound for every conceivable algorithm unless stated.

\subsection{Width and conditioning}
Theorem~\ref{thm:approx} is polynomial in $I$ and $\kappa$ for every fixed
$p$. Del Pia and Khajavirad proved that minimizing $x^\T Qx+c^\T x$ over
$[0,1]^n$ is strongly NP-hard when the interaction graph has treewidth two,
even with $\max\abs{Q_{ij}}\le5$ and $\max\abs{c_i}\le4$
\cite[Theorem~3]{DelPiaKhajavirad2026}. Hence, unless P${}={}$NP, no
polynomial-time algorithm exists for treewidth two without a further
parameter, and the instances produced by any such reduction must have
either several global minimizers or a condition number $\kappa$ that is
superpolynomial in $I$. 

The exponential dependence on $p$ cannot be removed either, even when
$L=0$.

\begin{proposition}\label{prop:eth}
Unless the exponential-time hypothesis fails, there is no algorithm that
solves every box QP $\min\{x^\T Qx+c^\T x:x\in[0,1]^n\}$ with zero diagonal,
given with a tree decomposition of bag size $p$, in time $2^{o(p)}\poly(I)$.
\end{proposition}

\begin{proof}
For a graph $G=(V,E)$, the zero-diagonal QP
$\min\{-\sum_ix_i+2\sum_{ij\in E}x_ix_j:x\in[0,1]^V\}$ has an optimal vertex
because $F$ is affine in each coordinate, and at a binary optimum the
selected vertices form an independent set (dropping an endpoint of a
selected edge changes the value by $1-2k\le-1$, where $k\ge1$ counts selected
neighbours). So the optimal value is minus the independence number. With the
single bag $V$, $p=\abs V$, and a $2^{o(p)}\poly(I)$ algorithm would compute
the independence number in time $2^{o(\abs V)}$, which contradicts the
exponential-time hypothesis \cite{ImpagliazzoPaturiZane2001}.
\end{proof}

Our algorithm solves this class in $O(2^p\poly(I))$ time (corrected grids with
$L=0$ consist of endpoints), so for $L=0$ the dependence on $p$ is optimal up
to the base of the exponent.

\subsection{Exact messages}
The exact alternative to grids is to represent each Bellman message as a
piecewise polynomial. On forests this works: the messages of a box QP are
concave piecewise quadratics with linearly many pieces
\cite{DelPiaKhajavirad2026}. At treewidth two it can fail even under the
assumptions of Theorem~\ref{thm:approx}.

\begin{proposition}[Expanding-box chain]\label{prop:chain}
For the function $G_m$ of Example~\ref{ex:chain}:
\begin{enumerate}[label=(\alph*)]
\item $\norm S^2+\norm z^2\le8\,G_m(S,z)$ on the box, so $0$ is the unique
minimizer, with growth constant $g=1/8$; the coordinate curvatures are
$10$ for $S_t$ ($t<m$), $4$ for $S_m$ and $7/4$ for $z_t$, so $\kappa=80$;
\item every representation of
$M(v)=\min\{G_m(S,z):S_m=v\}$ on $[0,2^m-1]$ by quadratic polynomials on
intervals uses at least $2^{m-1}$ pieces.
\end{enumerate}
\end{proposition}

\begin{proof}
(a) Write $r_t=S_t-2S_{t-1}-z_t$. Since $z_t(1-z_t)\ge0$,
$G_m\ge S_m^2+\norm r^2$. From $z_t\ge0$ and $S_{t-1}\ge0$,
$0\le S_{t-1}\le(S_t+\abs{r_t})/2$, and iterating,
$S_t\le2^{-(m-t)}S_m+\sum_{j>t}2^{-(j-t)}\abs{r_j}$. The vector
$(2^{-(m-t)})_t$ has squared norm at most $4/3$, and the operator
$\abs r\mapsto(\sum_{j>t}2^{-(j-t)}\abs{r_j})_t$ has norm at most
$\sum_{k\ge1}2^{-k}=1$, so $\norm S\le\frac2{\sqrt3}S_m+\norm r$ and
$\norm S^2\le\frac83S_m^2+2\norm r^2$. Also
$0\le z_t\le S_t+\abs{r_t}$, so $\norm z^2\le2\norm S^2+2\norm r^2$. Hence
$\norm S^2+\norm z^2\le3\norm S^2+2\norm r^2\le8S_m^2+8\norm r^2\le8G_m$.
The curvatures are the diagonal entries of the Hessian.
(b) Since $G_m\ge S_m^2$, $M(v)\ge v^2$. For an integer
$v=\sum_tz_t2^{m-t}$ with $z\in\{0,1\}^m$, the states
$S_t=\sum_{k\le t}z_k2^{t-k}\in[0,2^t-1]$ give $r=0$ and $G_m=v^2$. If
$M(v)=v^2$, the minimizer has $r=0$ and binary $z$, so $v$ is an integer.
Thus $M(v)-v^2$ vanishes exactly at the $2^m$ integers of $[0,2^m-1]$. On a
nondegenerate piece, $M(v)-v^2$ is a quadratic polynomial that is not
identically zero, so it has at most two zeros there; a degenerate piece
contains at most one. Hence at least $2^{m-1}$ pieces are needed.
\end{proof}

The variables $S_t$ have box widths $2^t-1$, so the input has $O(m^2)$ bits,
and the number of message pieces is exponential in the input length. After
the rescaling $s_t=2^{-t}S_t$ to unit boxes, the coordinate curvature becomes
$\Theta(4^m)$ while the growth constant stays $1/8$ ; the same
problem is then badly conditioned in our sense. The parameter $\kappa$
depends on the units, and good units can be essential.

\subsection{Local corrections and local moments}
Example~\ref{ex:star} shows that subtracting corrections only for the
coordinates of one bag, while the other coordinates are optimized on the
grid, does not give a valid bound. 

\subsection{Product domains}
Lemma~\ref{lem:interp} rounds coordinates independently. With the single
coupling constraint $x_1=x_2$ on $[0,1]^2$ and the grid $\{0,1\}$, the
feasible point $(\frac12,\frac12)$ is the average of no feasible grid
points other than $(0,0)$ and $(1,1)$, and the independent rounding puts
mass on the infeasible points $(0,1)$ and $(1,0)$. Correlated rounding
repairs this exactly when the fibers of the cells are integral, which is
the role of total unimodularity in Section~\ref{sec:constraints}. For a
general sparse linear or nonlinear constraint we know no rounding law with
a comparable error bound that preserves the decomposition.

\subsection{Several minimizers}
The localization argument needs a single minimizer. Let
$F(x)=2x_1^2+2x_2^2-5x_1x_2$ on $[-1,1]^2$. Its global minimizers are
$\pm(1,1)$ with value $-1$, and since $F$ is a quadratic, it has quadratic
growth towards the optimal set $\{\pm(1,1)\}$ \cite{LuoSturm2000}. Every
filtering step keeps both $-1$ and $1$ in each coordinate, so the coordinate
hulls never shrink, and a graded grid around one minimizer needs
$\Theta(\theta^{-1}\log(1/h))$ nodes per coordinate to reach mesh $h$ near the
center. The number of nodes then grows with the accuracy, and the
contraction argument fails because consecutive centers may approach
different minimizers. Keeping unions of retained cells instead of hulls
avoids the first problem when each coordinate has finitely many optimal
values (Section~\ref{sec:optsets}); for optimal sets with continuous
components we do not have an accuracy-independent bound.
